\chapter{Ejecucion Hibrida, Simulacion y Validacion Experimental}

\section{6.1 Introduccion: cierre del bucle planificacion-ejecucion}
Se formaliza la transicion desde solucion de optimizacion hacia comportamiento operacional verificable.

\section{6.2 Representacion formal del plan de ejecucion (Fase 2)}
Se define la estructura de plan, sus invariantes y su semantica de consumo por componentes runtime.

\section{6.3 Simulador hibrido: arquitectura y validacion topologica (Fase 2)}
Se describe el diseno del simulador y los chequeos de coherencia sobre el DAG de ejecucion.

\section{6.4 Estudio de prediccion: latencia, memoria y costos de transferencia (Fase 2)}
Se especifican metricas de prediccion y metodologia para evaluacion determinista previa a ejecucion real.

\section{6.5 Implementacion del ejecutor hibrido en PyTorch (Fase 3)}
Se documentara la implementacion operacional de la politica de asignacion en entrenamiento real.

\section{6.6 Baselines reales y comparativos experimentales (Fase 3)}
Se establecera comparacion sistematica frente a politicas de referencia en entorno fisico.

\section{6.7 Prediccion frente a observacion: contraste operacional (Fase 3)}
Se cuantificara el error de prediccion y su distribucion por regimen de carga.

\section{6.8 Exactitud final del modelo y estabilidad de convergencia (Fase 3-5)}
Se evaluara impacto de la ejecucion hibrida sobre calidad final del entrenamiento.

\section{6.9 Extensiones con persistencia de activaciones y scheduling asincrono (Fase 4)}
Se integraran extensiones avanzadas de memoria y planificacion para cerrar la propuesta completa.

\section{6.10 Analisis integrado de resultados y limites experimentales (Fase 5)}
Se sintetizaran resultados globales y restricciones metodologicas observadas en campana doctoral.
